%% Consignes : « Incontournable pour le rapport de MFE, l'etat de l'art
%% presente les solutions techniques, methodes, outils, organisations
%% existants qui repondent pour tout ou en partie a la problematique posee.
%% Il s'agit de realiser une analyse critique des travaux existants. »

The question of Section~\ref{sec:question} touches several fields that do not usually cite
one another. This section reviews each for what it settles and, more importantly, for what it
leaves open. A structured evaluation of every source, recording its argument, method and
limitations, was maintained throughout the mission and informs the appraisals below.

\subsection{Deadline-Driven Value and Risk Propagation}

Real-time systems research models the value of a task as a function of when it completes.
Chen and M\"uhlethaler show that utility typically holds constant until a deadline and then
decreases~\autocite{cheng1996due}, and a wider survey confirms structured degradation past
critical thresholds~\autocite{cheng2025survey}. This establishes that urgency is legitimately
time-varying.

Separately, supply chain research formalises the ripple effect, whereby a local disruption
cascades through downstream tiers and amplifies as it
travels~\autocite{ivanov2020digital,ivanov2019viable,li2020exploring}. This establishes that
urgency is legitimately network-dependent.

\emph{Critical appraisal.} The first stream stops at a single isolated task and offers
nothing on combining heterogeneous indicators. The second models propagation but treats the
node-level score as given. Neither addresses a declared signal at all. Between them they
justify treating urgency as time-varying and propagated, and leave the aggregation question
entirely open.

\subsection{Indicators of Delivery Rupture}

If urgency is to be aggregated, the indicators must come from somewhere. Heckmann, Comes and
Nickel establish rupture as the terminal event a supply risk measure should anticipate, and
criticise the field for measuring exposure without connecting it to that
event~\autocite{heckmann2015critical}. Indicator-based supplier monitoring proceeds by
identifying observable precursors~\autocite{mokhtar2019supplier}.

Five families recur. Slack is the only indicator whose exhaustion is
irreversible~\autocite{cohen2021assigning,sun2020combating}. Capacity saturation adds waiting
time a nominal lead time does not
show~\autocite{tomlin2006value,hosseini2019resilient}. Resource performance is monitored
through composite proxies because supplier-side data is rarely
available~\autocite{mokhtar2019supplier,hlioui2017joint}. Adverse-event exposure is treated
probabilistically~\autocite{hosseini2020ripple,liu2021robust}. Cost signals are read as
symptoms of strain rather than causes~\autocite{shen2019managing}, and emissions constraints
restrict the admissible responses rather than causing
rupture~\autocite{hoen2014effect,smartfreight2023glec}.

\emph{Critical appraisal.} This literature names precursors and monitors them individually.
It supplies no criterion for deciding which belong in one composite, and says nothing about
whether equal numerical values on structurally different indicators carry equal risk meaning.
Both questions are inherited by any multi-block model, and
Section~\ref{sec:realisation_modele} treats the second explicitly.

\subsection{Cognitive Bias in Declared Urgency}

Zhu, Yang and Hsee report the mere-urgency effect across five controlled experiments:
participants consistently preferred imminent low-value tasks to distant high-value ones, and
the bias survived being told about it~\autocite{zhu2018mere}. Steel gives a formal
discounting model in which immediacy becomes more salient as a deadline
approaches~\autocite{steel2006integrating}. Field studies confirm the pattern under
workload~\autocite{rusou2020psychology,diwa2020task}.

\emph{Critical appraisal.} This stream diagnoses the problem and quantifies it in controlled
settings, but the magnitude in industrial supply chains is uncalibrated, and it prescribes no
operational correction. It justifies treating the declared signal as biased; it does not say
what to do about it.

\subsection{Multi-Criteria Elicitation and Aggregation}

Structured elicitation of a human judgment is a solved problem in outline. The Analytic
Hierarchy Process derives criteria weights from pairwise comparisons and validates them
through a consistency ratio~\autocite{saaty1980analytic}, though it is vulnerable to rank
reversal and to cognitive noise~\autocite{dyer1990,triantaphyllou2024}. The Best-Worst Method
reduces the number of comparisons required~\autocite{rezaei2015best}, and its fuzzy extension
absorbs judgment uncertainty~\autocite{guo2017fuzzy,liang2020}. For ranking alternatives
without letting a good score on one criterion mask a catastrophic one on another, outranking
methods such as PROMETHEE~II establish pairwise preference
flows~\autocite{brans1986promethee,behzadian2010promethee}.

The composite indicator literature supplies the necessary warning: nominal weights often
diverge from effective importance because of correlation and normalisation
effects~\autocite{greco2018,becker2017,borgonovo2021global}.

\emph{Critical appraisal.} These are static tools. They elicit and aggregate, and none is
routinely compared against a computed operational counterpart. The composite-index critique
is generally raised as a caveat at the end of a paper rather than turned into something a
running system checks on its own data.

\subsection{Digital Twins and Decision Support}

The digital twin concept was formalised by Grieves~\autocite{grieves2014digital} and extended
to manufacturing~\autocite{tao2018digital,fuller2020digital}. In supply chains it integrates
simulation, optimisation and machine learning~\autocite{ivanov2021digital}, with demonstrated
capability in disruption monitoring and
replanning~\autocite{ding2023realtime,leung2022digital,hrusovsky2021realtime}. Recent
human-in-the-loop cognitive twins capture human feedback to synchronise mental models with
simulated operations~\autocite{niloofar2023}.

\emph{Critical appraisal.} This is the family the delivered system belongs to, and the gap is
here. No retrieved system combines a declared urgency input, a computed operational urgency
and their propagated mismatch. A second gap matters for this mission specifically: these
systems are evaluated on how accurately they estimate physical state, and their inputs are
assumed to arrive. Where an input depends on the same operator whose judgment the system
exists to support, that circularity is not treated as a measurement problem.

\subsection{Verifying a Probabilistic Signal}

Because this mission scores a probabilistic signal against outcomes, it inherits forecast
verification. Brier introduced the mean squared error of probabilistic
forecasts~\autocite{brier1950verification}, and Murphy decomposed it into reliability,
resolution and uncertainty~\autocite{murphy1973vector}, establishing the distinction the
results of this report depend on: a forecast can order outcomes correctly while announcing
the wrong probabilities, and the two failures need different repairs. Machine learning
rediscovered the same point, finding strong classifiers that produce distorted probabilities
correctable by a cheap post-hoc layer~\autocite{niculescu2005predicting,guo2017calibration}.

\emph{Critical appraisal.} This literature assumes the outcome being predicted is
unambiguously defined. In an operational supply chain it is a modelling choice, and
Section~\ref{sec:realisation_resultats} measures that choice moving the reported result by a
factor of five.

\subsection{What the Review Establishes}
\label{sec:gap}

Each stream closes on a limit. Collected, those limits describe one absence rather than
several. The reviewed work covers multi-block risk aggregation, network propagation, digital
twin operation, multi-criteria elicitation and probabilistic verification. No retrieved work
combines them with a structured declared-urgency input and an empirical test of the resulting
adequacy signal against recorded outcomes.

That missing combination is what this mission set out to build and to test.
